\section{Averaging-Based Removal of Shared Linguistic Structure}
\label{sec:avg_remove_linguistic}

We formalize and analyze an averaging-based approximation method for removing
shared linguistic structure from high-dimensional token representations.
The goal is to show that simple averaging removes the dominant, anisotropic
linguistic component while preserving semantic variability, thereby providing
a hardware-friendly alternative to SVD-based spike removal.

\subsection{Model Assumptions}

Let the ambient feature space be $\mathbb{R}^D$, decomposed as
\[
\mathbb{R}^D = \mathcal{L} \oplus \mathcal{S},
\]
where:
\begin{itemize}
    \item $\mathcal{L}$ is a low-dimensional \emph{linguistic subspace}
          of dimension $d_\ell \ll D$, capturing shared grammatical and structural
          regularities common across tokens;
    \item $\mathcal{S}$ is a high-dimensional \emph{semantic subspace}
          of dimension $d_s = D - d_\ell$, capturing token-specific meaning.
\end{itemize}

Let $\{x_i\}_{i=1}^N \subset \mathbb{R}^D$ denote token representations in a batch.
Each token is assumed to decompose as
\begin{equation}
\label{eq:token_decomposition}
x_i = \ell_i + s_i,
\end{equation}
where:
\begin{itemize}
    \item $\ell_i \in \mathcal{L}$ is a \emph{linguistic component}.
          While $\ell_i$ may vary across sequences or contexts, all $\ell_i$
          lie in the same low-dimensional subspace $\mathcal{L}$;
    \item $s_i \in \mathcal{S}$ is a \emph{semantic component}.
\end{itemize}

This formulation explicitly allows linguistic structure to exhibit
\emph{controlled variance} across sequences, while remaining confined to a
low-rank subspace.

\paragraph{Assumption 1 (Semantic compressibility).}
Each $s_i$ lies in $\mathcal{S}$ and is \emph{compressible} in the sense that
its energy is concentrated on a subset of coordinates whose effective dimension
is much smaller than $d_s$.
Equivalently, for any fixed energy fraction $\alpha \in (0,1)$, the number of
coordinates required to capture $\alpha\|s_i\|_2^2$ satisfies
\[
k_\alpha(s_i) \ll d_s.
\]
This reflects that each token activates only a limited number of semantic
features, even though the semantic space itself is large.

\paragraph{Assumption 2 (Semantic diversity under batch sampling).}
Tokens in a batch are drawn from different sentences or sufficiently distant
contexts such that their semantic components $\{s_i\}$ can be treated as
approximately independent samples from a common distribution on $\mathcal{S}$
with finite second moment.
Consequently, the empirical semantic mean concentrates:
\begin{equation}
\label{eq:semantic_mean_concentration}
\left\|
\frac{1}{N}\sum_{i=1}^N (s_i - \mathbb{E}s_i)
\right\|
= O(N^{-1/2})
\quad \text{with high probability}.
\end{equation}

Importantly, semantic components need not have zero mean exactly; it suffices
that any batch-level semantic bias is small compared to the shared linguistic
structure and vanishes under aggregation.

These assumptions are consistent with empirical observations in large language
models: linguistic structure is coherent and low-rank, while semantic information
is high-dimensional, diverse, and weakly correlated across tokens.

\subsection{Second-Moment Matrices}

Define the uncentered second-moment matrix
\begin{equation}
\label{eq:M_def}
M \triangleq \frac{1}{N}\sum_{i=1}^N x_i x_i^\top,
\end{equation}
and the centered covariance matrix
\begin{equation}
\label{eq:C_def}
C \triangleq \frac{1}{N}\sum_{i=1}^N (x_i - \mu)(x_i - \mu)^\top,
\qquad
\mu \triangleq \frac{1}{N}\sum_{i=1}^N x_i.
\end{equation}

\subsection{Effect of Averaging}

Using \eqref{eq:token_decomposition}, we compute the sample mean:
\[
\mu
= \frac{1}{N}\sum_{i=1}^N \ell_i
  + \frac{1}{N}\sum_{i=1}^N s_i.
\]
By Assumption~2, the semantic mean term concentrates and is $O(N^{-1/2})$.
We therefore decompose
\[
\mu = \bar{\ell} + O(N^{-1/2}),
\qquad
\bar{\ell} \triangleq \frac{1}{N}\sum_{i=1}^N \ell_i \in \mathcal{L}.
\]

Substituting into \eqref{eq:M_def}, we obtain
\begin{align}
M
&= \frac{1}{N}\sum_{i=1}^N (\ell_i + s_i)(\ell_i + s_i)^\top \\
&= \underbrace{\frac{1}{N}\sum_{i=1}^N \ell_i \ell_i^\top}_{\text{linguistic spike } L}
 + \underbrace{\frac{1}{N}\sum_{i=1}^N s_i s_i^\top}_{\text{semantic covariance } S}
 + \text{cross terms}.
\end{align}
The cross terms vanish up to $O(N^{-1/2})$ by Assumption~2.
Hence,
\begin{equation}
\label{eq:M_decomp}
M = L + S + O(N^{-1/2}),
\qquad
\mathrm{rank}(L) \le d_\ell.
\end{equation}

Similarly, the centered covariance satisfies the exact identity
\begin{equation}
\label{eq:C_M_relation}
C = M - \mu\mu^\top
  = S + O(N^{-1/2}).
\end{equation}

\subsection{Main Theorem}

\begin{theorem}[Averaging removes shared linguistic structure while preserving semantics]
\label{thm:avg_removes_linguistic}
Under the assumptions above:
\begin{enumerate}
    \item The uncentered second-moment matrix $M$ decomposes into a low-rank
          linguistic spike $L$ with $\mathrm{rank}(L)\le d_\ell$ plus a semantic
          covariance $S$, up to a vanishing $O(N^{-1/2})$ term.
    \item The centered covariance matrix $C$ equals the semantic covariance $S$
          up to the same vanishing term.
    \item For any vector $v \in \mathbb{R}^D$ such that $v \perp \mathcal{L}$,
          \[
          Mv = Cv + O(N^{-1/2}).
          \]
          Hence, $M$ and $C$ induce asymptotically identical linear operators on
          $\mathcal{L}^\perp$.
\end{enumerate}
Consequently, subtracting the mean via averaging removes only the shared
linguistic subspace and preserves all semantic directions relevant for learning.
\end{theorem}

\begin{proof}
The decomposition \eqref{eq:M_decomp} follows from bilinearity and the
concentration of semantic averages.
Equation~\eqref{eq:C_M_relation} is an identity.
For $v \perp \mathcal{L}$, we have $v^\top \ell_i = 0$ for all $i$, and therefore
\[
Mv
= (L + S)v + O(N^{-1/2})
= Sv + O(N^{-1/2})
= Cv + O(N^{-1/2}),
\]
which proves the claim.
\end{proof}

\subsection{Implications for Approximation Algorithms}

Theorem~\ref{thm:avg_removes_linguistic} shows that subtracting the mean
(i.e., simple averaging) removes the dominant low-rank linguistic structure
while leaving semantic covariance---the second-order structure that drives
learning and generalization---essentially unchanged.

In contrast, SVD-based methods (e.g., Metis/Matisse) explicitly compute and remove
the top singular subspace associated with $\mathcal{L}$, which is computationally
and hardware expensive.

Averaging provides a principled approximation to this operation with:
\begin{itemize}
    \item $O(ND)$ cost (single pass over tokens),
    \item no matrix decomposition,
    \item preservation of semantic information up to vanishing error.
\end{itemize}

Thus, averaging-based spike removal offers an efficient and theoretically grounded
alternative to SVD-based linguistic spike suppression in high-dimensional token
representations.